\documentclass[10pt,a4paper]{amsart}
\usepackage[margin=19mm]{geometry}
\usepackage{amsmath,amssymb,mathtools}
\usepackage[scr=rsfs,cal=euler]{mathalfa}
\usepackage{xfrac}
\usepackage[unicode,hidelinks,bookmarks=false]{hyperref}
\newcommand{\ie}{i.e.\ }
\newcommand{\cf}{cf.\ }
\newcommand{\resp}{respectively\ }
\newcommand{\Subsection}{\subsection}
\providecommand{\nobreakdash}{\nobreak\hbox{-}}
\setlength{\emergencystretch}{2em}
\multlinegap0pt
\usepackage{amssymb}
\usepackage{xcolor}
\usepackage{algorithm}\usepackage{algpseudocode}
\newtheorem{lemma}{Lemma}
\title{Dynamical Systems and Nonconvex Quadratic Programs}
\author{Nguyen Nang Thieu}
\date{Source: arXiv:2609.06088v1 (CC BY 4.0)}
\begin{document}
\maketitle

\noindent\emph{Source and licence.} Dynamical Systems and Nonconvex Quadratic Programs. Version arXiv:2609.06088v1 (CC BY 4.0), distributed by arXiv under the Creative Commons Attribution 4.0 International licence, http://creativecommons.org/licenses/by/4.0/, as recorded in the arXiv licence metadata for this exact version and shown on the abstract page https://arxiv.org/abs/2609.06088v1. Full licence notice: \texttt{licenses/arxiv-2609.06088-CC-BY-4.0.txt}.\par\smallskip

\noindent\textit{Editorial scope.} Counted unit: the article's lemma lem4 (source0.tex lines 550--559). The article's complete original proof of it is delivered with it (source0.tex lines 560--582, one \texttt{proof} environment of 23 lines). No mathematics is written, repaired or re-derived: the statement and the proof are the article's own lines, in the article's own notation, and no retained line is substituted or re-flowed.  Retained uncounted as the source-local context the proof needs: lines 503--504.  Nothing was omitted from any retained span, so the package carries no omission record.  No retained line contains a citation, so the wrapper carries no bibliography and no bibliography entry.  No retained line contains a figure environment, an image-inclusion command or drawing code, so no image is delivered and the package declares no asset dependency.  Editorial formatting only, in addition: an AMS \texttt{amsart} preamble that restates the article's own declarations and macros used by the retained lines, namely the environments \texttt{lemma} on separate counters, together with the standard packages those lines need.  The retained environments print in this document as Lemma 1 in order of appearance, whereas the article prints the same items with the numbering of its own full document.  The complete native source of the article as distributed in the arXiv e-print for this exact version is bundled unchanged as \texttt{sources/209491/source0.tex}.
where \begin{equation}\label{alpha}\alpha := \max\limits_{\lambda \in \sigma(Q)} \left[ \frac{\lambda}{\rho} \left( \frac{\lambda}{4\rho} - 1 \right)\right].
\end{equation}

\begin{lemma}\label{lem4}
	Let $\rho>0$ and $\alpha$ be given by~\eqref{alpha}. Then,
$$
	\begin{cases}
		\alpha <0, & \text{if } \sigma(Q)\subset(0,4\rho),\\[2mm]
		\alpha =0, & \text{if } \sigma(Q)\subset[0,4\rho] \text{ and } \sigma(Q)\cap\{0,4\rho\}\neq\emptyset,\\[2mm]
		\alpha >0, & \text{if } \sigma(Q)\not\subset[0,4\rho].
	\end{cases}
$$
\end{lemma}

\begin{proof}
	Put
	$\varphi(\lambda)=\frac{\lambda^2}{4\rho^2}-\frac{\lambda}{\rho}$ for $\lambda\in\mathbb{R}$, and note by~\eqref{alpha} that  $\alpha=\max\limits_{\lambda\in\sigma(Q)}\varphi(\lambda)$. Since $0$ and~$4\rho$ are the two roots of the quadratic polynomial $\varphi(\lambda)$ whose highest coefficient is positive, we have
	\begin{equation}\label{value_varphi}
		\begin{cases}
			\varphi(\lambda)<0 \ \text{ for } \lambda\in(0,4\rho), \\
			\varphi(0)=\varphi(4\rho)=0, \\
			\varphi(\lambda)>0 \ \text{ for } \lambda\in(-\infty,0)\cup(4\rho,\infty).
		\end{cases}
	\end{equation}
	
If $\sigma(Q)\subset(0,4\rho)$, then by~\eqref{value_varphi} one has $\alpha=\max\limits_{\lambda\in\sigma(Q)}\varphi(\lambda)<0$.

If $\sigma(Q)\subset[0,4\rho]$ and $\sigma(Q)\cap\{0,4\rho\}\neq\emptyset$ then, using~\eqref{value_varphi} again, one can deduce that~$\alpha=0$. 

If $\sigma(Q)\not\subset[0,4\rho]$, then there exists 
 $\bar\lambda\in\sigma(Q)$ with $\bar\lambda<0$ or there is  $\hat\lambda\in\sigma(Q)$ with $\hat\lambda>4\rho$. Since $\varphi(\bar\lambda)>0$ and $\varphi(\hat\lambda)>0$ by~\eqref{value_varphi}, one has
$
\alpha=\max\limits_{\lambda\in\sigma(Q)}\varphi(\lambda)>0.
$

The proof is complete.
\end{proof}

\end{document}
